\documentclass[journal]{IEEEtran}
\usepackage{amsmath,amssymb,amsthm,graphicx,booktabs,array,cite,url}
\usepackage[hidelinks]{hyperref}
\newtheorem{theorem}{Theorem}
\newtheorem{lemma}{Lemma}
\newtheorem{proposition}{Proposition}
\newtheorem{corollary}{Corollary}
\newcommand{\norm}[1]{\left\lVert#1\right\rVert}
\newcommand{\ket}[1]{\lvert#1\rangle}
\newcommand{\bra}[1]{\langle#1\rvert}
\newcommand{\R}{\mathbb{R}}
\input{generated/metrics}
\begin{document}
\title{A Resource-Explicit Chebyshev-LCU Baseline for Coherent Matrix Inversion: Theory and Emulation to Order 1000}
\author{Vahid Tavakkoli, Kabeh Mohsenzadegan, and Kyandoghere Kyamakya%
\thanks{QCeNN-MI research project. Reproducible implementation and experiment sources: \url{https://github.com/vtavakkoli/QCeNN-MI}. This manuscript reports simulator evidence; physical quantum-device execution is future work.}}
\maketitle
\begin{abstract}
Quantum inverse-action circuits can encode an $N$-dimensional vector with
$\lceil\log_2 N\rceil$ data qubits, but this representation does not remove
matrix-loading, conditioning, gate-depth, or readout costs. This paper develops
a resource-explicit spectral baseline for real nonsingular matrices. An analytic
Chebyshev reciprocal expansion is implemented by linear combination of unitaries
(LCU), a single walk ancilla, and a binary selector. We derive a uniform operator
residual bound, relative solution and normalized-state bounds, and a conditional
success-probability guarantee. The construction uses classical dense spectral or
singular-value decomposition; it is therefore not an efficient black-box input
oracle or a quantum-speedup claim. A structured gate-statevector simulator
retains all selector and walk amplitudes and evaluates selector uncomputation
through an exact success projection. The reproducible validation package checks
\TheoryCaseCount{} degree/conditioning cases and audits the complete success
operator for \PrecisionCaseCount{} real $1000\times1000$ matrices. The recorded
largest operator residual is \WorstOperatorResidual{} with a polynomial target
of $10^{-12}$ and a separate numerical acceptance limit of $10^{-10}$.
The large instances use 15--21 raw qubits. At condition number 1000, the
worst-input success probability is approximately $10^{-6}$. Native-gate and
synthetic-noise checks through dimension 100 expose the distinction between
ideal numerical accuracy and useful hardware quality. A single containerized
workflow regenerates numerical evidence, resource logs, tables, figures, and
an accompanying hardware-access assessment report.
\end{abstract}
\begin{IEEEkeywords}
Quantum linear systems, coherent inverse action, Chebyshev approximation,
linear combination of unitaries, reproducible emulation, resource estimation.
\end{IEEEkeywords}

\section{Introduction}
\IEEEPARstart{M}{atrix} inversion is often posed as the search for $X$ satisfying
$MX=I$. In applications, the desired object may instead be one inverse action
$x=M^{-1}b$, a normalized quantum state proportional to that vector, or an
observable of the solution. These outputs are not interchangeable. A quantum
state may represent a long vector compactly without making its classical
entries inexpensive to recover. The distinction is central to quantum
linear-system algorithms and to credible evaluations of their implementations
\cite{harrow2009,childs2017,gilyen2019}.

The QCeNN-MI project explores matrix-inversion dynamics motivated by recurrent
neural formulations of time-varying inversion \cite{tavakkoli2019}. Its early
experiments comprise small variational cells, coupled ansatz templates,
stateful classical tracking, and safeguarded polynomial corrections. Those
hybrid experiments motivate a separate question: can an input-independent
coherent operation approximate inverse action without fitting the solution
state or applying a classical residual correction inside the solver?

This paper studies that question through a deliberately explicit compiled
baseline. A dense classical decomposition supplies input singular bases and
positive singular values. A reciprocal polynomial, whose coefficients depend
only on the extremal singular values, is implemented in a unitary dilation.
Ancilla postselection exposes the inverse-action block. The method is
standalone at the level of the coherent inverse-action subroutine, while the
encoding remains classical and dense. It is not a physical realization of a
local quantum cellular network, nor a conversion of every hybrid tracking
step into a quantum operation.

The contributions are fourfold. First, we give a self-contained certification
of the particular single-flag Chebyshev-LCU construction used in the code,
including its residual, output-state, and success bounds. The approximation
and LCU ingredients are established techniques, not new general quantum
linear-system theory. Second, we expose the conditions under which the
logarithmic data-register count coexists with substantial encoding and
sampling costs. Third, we audit full simulated operators at dimension 1000
and cross-check manageable instances with independently executed native
circuits. Fourth, we supply an executable evidence pipeline that retains
failed noisy-quality outcomes and supports a staged physical-device study.

Our main empirical result is ideal-simulation feasibility for the tested
families with quantified error. It does not prove uniform performance over
all nonsingular matrices. The exact-arithmetic convergence theorem applies
to a finite positive singular-value interval; degree, floating-point
resolution, and postselection resources remain explicit parameters.

\section{Problem, Encoding, and Evidence Model}
\subsection{Domain and outputs}
Let $M\in\R^{N\times N}$ be nonsingular, with singular-value decomposition
\begin{equation}
 M=U\Sigma V^T,\qquad
 \Sigma=\operatorname{diag}(\sigma_1,\ldots,\sigma_N),
 \quad 0<a\leq\sigma_j\leq b.
 \label{eq:svd}
\end{equation}
Write $\kappa=b/a$. For a nonzero real right-hand side $b_0$, define
$\ket{b_0}=b_0/\norm{b_0}_2$ and the desired state
$\ket{x}=M^{-1}b_0/\norm{M^{-1}b_0}_2$. The name $b_0$ separates the input
vector from the upper singular bound $b$. The implementation currently
rejects complex, non-square, singular, and numerically singular inputs.
Rectangular pseudoinversion is outside the theorem and the large-scale study.

For real square $M$,
\begin{equation}
 MM^T\succ0\quad\Longleftrightarrow\quad M\text{ is nonsingular}.
 \label{eq:invertible}
\end{equation}
Indeed, $y^TMM^Ty=\norm{M^Ty}_2^2$ is strictly positive for every nonzero
$y$ exactly when $M^T$ has a trivial kernel. Positive definiteness is not
elementwise positivity. Condition~\eqref{eq:invertible} asserts existence of
an inverse; it places no useful uniform upper bound on $\kappa$. The solver
does not explicitly form or solve the normal equations $MM^T$, which would
square the spectral condition number.

An arbitrary inverse is generally nonunitary. Consequently, the target of
this work is a postselected block of a larger unitary. The successful circuit
output is a normalized solution state; the matrix norm scale remains encoded
in the known LCU normalization and success probability. Full classical
inverse reconstruction is a separate measurement and output task.

\subsection{Registers and padding}
Set $n=\lceil\log_2N\rceil$ and $P=2^n$. Valid amplitudes occupy the first
$N$ data coordinates. The singular bases are embedded as block-diagonal
orthogonal transformations on the $P$-dimensional data register. Padded
coordinates are not populated by the prepared input and are not coupled to
valid coordinates. Thus non-power-of-two dimensions do not require one qubit
per row, but padding can enlarge a dense compilation by up to a constant factor.

For a degree-$d$ polynomial, $s=\lceil\log_2(d+1)\rceil$ selector qubits and
one walk ancilla are used, giving
\begin{equation}
 q_{\mathrm{total}}=\lceil\log_2N\rceil+1+
                    \lceil\log_2(d+1)\rceil.
 \label{eq:qubits}
\end{equation}
These are raw algorithmic qubits. Provider compiler ancillas and quantum
error-correction overhead are excluded.

\subsection{What is classical and what is coherent}
The decomposition~\eqref{eq:svd}, scalar interval analysis, coefficient
construction, and circuit synthesis are classical preprocessing. For SPD
matrices a positive eigendecomposition is used; general real matrices use
SVD. A known input may be prepared directly in the $U^T$ basis to reduce gate
count. An unfused input-independent operator is also exported and tested on
different previously prepared inputs.

The coherent core comprises selector preparation, controlled walk powers,
coefficient signs, selector unpreparation, and the output basis change.
No classical inverse, reference solution, inverse-eigenvalue table, fitted
solution state, or residual-feedback update is supplied to that core. An
independent classical solve is performed afterwards for audit or reference
projection only. Since the classical decomposition already incurs dense
cubic work, this separation does not establish end-to-end quantum advantage.

We distinguish four evidence levels: analytic exact-arithmetic guarantees;
structured ideal gate-statevector emulation; native-gate ideal and synthetic
noise emulation; and physical-device measurements. Only the first three are
present. Success at one level must not be relabeled as evidence at another.

\section{Reciprocal Approximation and Certification}
\subsection{An analytic Chebyshev series}
Assume first that $a<b$, and define
\begin{equation}
 \begin{gathered}
 c=\frac{a+b}{2},\quad h=\frac{b-a}{2},\quad t=\frac{\sigma-c}{h},\\
 r=\frac{\sqrt b-\sqrt a}{\sqrt b+\sqrt a}\in(0,1).
 \end{gathered}
 \label{eq:parameters}
\end{equation}
With $T_k(\cos\theta)=\cos(k\theta)$, let
\begin{equation}
 p_d(\sigma)=\frac1{\sqrt{ab}}
 \left[1+2\sum_{k=1}^{d}(-r)^kT_k(t)\right].
 \label{eq:poly}
\end{equation}
This interval reciprocal is the particular Chebyshev representation used
here; Chebyshev-series operator implementations for linear systems are
established in the LCU literature \cite{childs2017}.

\begin{lemma}[Uniform truncation error]
For $\sigma\in[a,b]$,
\begin{equation}
 \left|p_d(\sigma)-\sigma^{-1}\right|
 \leq E_d:=\frac{2}{\sqrt{ab}}\frac{r^{d+1}}{1-r}.
 \label{eq:tail}
\end{equation}
\end{lemma}
\begin{proof}
The geometric-series identity yields the Poisson kernel
\[
 1+2\sum_{k=1}^{\infty}(-r)^k\cos(k\theta)
 =\frac{1-r^2}{1+2r\cos\theta+r^2}.
\]
The definitions imply $c/h=(1+r^2)/(2r)$ and
$\sqrt{ab}/h=(1-r^2)/(2r)$. Multiplying the kernel by $1/\sqrt{ab}$ therefore
gives $(c+h\cos\theta)^{-1}$. Since $|T_k(t)|\leq1$ on $[-1,1]$, the omitted
terms are bounded by the absolute geometric tail, giving~\eqref{eq:tail}.
\end{proof}

Define the inverse approximation
\begin{equation}
 X_d=V p_d(\Sigma)U^T,\qquad
 \delta_d=bE_d=2\sqrt\kappa\frac{r^{d+1}}{1-r}.
 \label{eq:delta}
\end{equation}

\begin{theorem}[Operator and per-input certification]
Under~\eqref{eq:svd},
\begin{align}
 \norm{MX_d-I}_2&\leq\delta_d, &
 \norm{X_dM-I}_2&\leq\delta_d,\label{eq:operatorbound}\\
 \norm{X_d-M^{-1}}_2&\leq E_d. &&\label{eq:inversebound}
\end{align}
For every nonzero $b_0$, writing $x_d=X_db_0$ and $x=M^{-1}b_0$ gives
\begin{equation}
 \frac{\norm{Mx_d-b_0}_2}{\norm{b_0}_2}\leq\delta_d,
 \qquad
 \frac{\norm{x_d-x}_2}{\norm{x}_2}\leq\delta_d.
 \label{eq:perinput}
\end{equation}
If $\delta_d<1$, then $x_d\ne0$ and
\begin{equation}
 \left\lVert\frac{x_d}{\norm{x_d}_2}
             -\frac{x}{\norm{x}_2}\right\rVert_2\leq2\delta_d.
 \label{eq:statebound}
\end{equation}
\end{theorem}
\begin{proof}
The two residuals are orthogonal similarities of
$\operatorname{diag}(\sigma_jp_d(\sigma_j)-1)$. Their spectral norms are at
most $bE_d$. Similarly, $X_d-M^{-1}$ has singular values
$|p_d(\sigma_j)-1/\sigma_j|$, proving~\eqref{eq:inversebound}.
The first inequality in~\eqref{eq:perinput} is the induced operator bound.
For the second, let $z=U^Tb_0$ and
$e_j=\sigma_jp_d(\sigma_j)-1$. Then
\[
 \norm{x_d-x}_2^2
 =\sum_j |e_j|^2|z_j/\sigma_j|^2
 \leq\delta_d^2\norm{x}_2^2.
\]
If $\delta_d<1$, this prevents $x_d=0$. For any two nonzero vectors $u,v$,
the triangle inequality and reverse triangle inequality imply
\[
 \norm{u/\norm{u}_2-v/\norm{v}_2}_2
 \leq 2\norm{u-v}_2/\norm{v}_2.
\]
Substitute $u=x_d$, $v=x$ to obtain~\eqref{eq:statebound}.
\end{proof}

\begin{corollary}[Finite-degree convergence]
For every fixed nonsingular real square matrix and every $\epsilon>0$, a
finite degree attains $\delta_d\leq\epsilon$ in exact arithmetic. A sufficient
choice for $a<b$ is
\begin{equation}
 d\geq
 \left\lceil\frac{\log(2\sqrt\kappa/((1-r)\epsilon))}{-\log r}\right\rceil-1.
 \label{eq:degree}
\end{equation}
If $a=b$, the reciprocal is constant and no nonconstant term is needed.
\end{corollary}
For large $\kappa$, $-\log r$ is asymptotic to $2/\sqrt\kappa$, giving the
sufficient scaling $O(\sqrt\kappa[\log\kappa+\log(1/\epsilon)])$. This is a
polynomial-degree statement for a known spectral interval, not an oracle
query-complexity result. No fixed degree budget works uniformly as the
smallest singular value tends to zero.

\section{Single-Flag Coherent LCU Construction}
\subsection{Walk and coefficient selection}
For $t_j=(\sigma_j-c)/h$ and $\theta_j=\arccos(t_j)$, define an encoded walk
whose flag block on spectral coordinate $j$ is
\begin{equation}
 R_j=\begin{bmatrix}
 \cos\theta_j&\sin\theta_j\\-\sin\theta_j&\cos\theta_j
 \end{bmatrix}.
 \label{eq:walk}
\end{equation}
It is a real unitary rotation, and
$\bra0R_j^k\ket0=\cos(k\theta_j)=T_k(t_j)$.
Let $c_0=1/\sqrt{ab}$ and $c_k=2(-r)^k/\sqrt{ab}$ for $k\geq1$, and set
\begin{equation}
 \alpha_d=\sum_{k=0}^d|c_k|,
 \quad \ket w=\sum_{k=0}^d\sqrt{|c_k|/\alpha_d}\ket k.
 \label{eq:lcuweights}
\end{equation}
Unused binary selector states have zero amplitude. If PREP maps $\ket0$ to
$\ket w$, SELECT applies $(-1)^kR^k$ on branch $k$. The sign is a single $Z$
on the least significant selector bit. The successful selector-and-flag block
of PREP$^\dagger$ SELECT PREP is
\begin{equation}
 (\bra{0^s}\bra0)\mathcal U_d(\ket{0^s}\ket0)
 =p_d(\Sigma)/\alpha_d.
 \label{eq:block}
\end{equation}
Applying $U^T$ before this block and $V$ afterwards gives $X_d/\alpha_d$ on
the valid data subspace. The operation is independent of the unknown solution
and, in its unfused form, independent of the chosen input vector.

\subsection{Binary powers and their access assumption}
For selector bits $k=\sum_{\ell=0}^{s-1}2^\ell k_\ell$, commuting rotations
give $R_j^k=\prod_\ell R_j^{2^\ell k_\ell}$. Because every $\theta_j$ is
classically known in this baseline, each controlled power is synthesized by
a uniformly controlled rotation with angle $-2^{\ell+1}\theta_j$.
Thus the exported circuit has one multiplexed walk-power gadget per selector
bit, rather than one fresh flag per polynomial factor.

This convenience is conditional on explicit spectral access. It does not
say that an arbitrary black-box walk can be powered at unit cost, nor that
angle synthesis is independent of required precision. Generic block-encoding
and QSVT results address different access models \cite{gilyen2019}; replacing
our classical spectral decomposition by such an oracle is future work.

\begin{proposition}[Normalization and success]
For $a<b$,
\begin{equation}
 \alpha_d=\frac1{\sqrt{ab}}
 \left(1+\frac{2r(1-r^d)}{1-r}\right)\leq\frac1a.
 \label{eq:alpha}
\end{equation}
For normalized input $\ket{b_0}$, postselection succeeds with probability
\begin{equation}
 p_{\mathrm{succ}}=\norm{X_d\ket{b_0}}_2^2/\alpha_d^2.
 \label{eq:success}
\end{equation}
If $\delta_d<1$, then for every input
\begin{equation}
 p_{\mathrm{succ}}\geq
 \frac{(1-\delta_d)^2}{\alpha_d^2b^2}
 \geq\frac{(1-\delta_d)^2}{\kappa^2}.
 \label{eq:successbound}
\end{equation}
\end{proposition}
\begin{proof}
Summing the absolute geometric coefficients gives~\eqref{eq:alpha}.
The infinite sum equals $a^{-1}$ because $(1+r)/(1-r)=\sqrt\kappa$.
Equation~\eqref{eq:block} and orthogonality of the input/output basis changes
give~\eqref{eq:success}. Since
$|\sigma_jp_d(\sigma_j)-1|\leq\delta_d<1$, each
$p_d(\sigma_j)\geq(1-\delta_d)/b$. This lower-bounds the smallest singular
value of $X_d$, proving~\eqref{eq:successbound}.
\end{proof}

The bound explains the approximately $\kappa^{-2}$ worst-input success in
the experiment. The absence of a large qubit register does not eliminate
sampling cost. Straight independent repetition can require order
$\kappa^2$ attempts per accepted output in the worst case. Amplitude
amplification is not implemented here, and its possible use must include
the cost of controlled reflections and repeated data access.

\section{Numerical Stability, Simulation, and Audit}
\subsection{Stable inputs and degree limits}
The implementation preserves the original decomposition singular values.
Recovering $\sigma$ by reversing an intermediate contraction
$1-\tau\sigma$ unnecessarily loses precision when subtraction cancels.
Walk angles are instead formed with the equivalent endpoint-stable expression
\begin{equation}
 \theta_j=2\operatorname{atan2}
       \left(\sqrt{b-\sigma_j},\sqrt{\sigma_j-a}\right).
 \label{eq:stableangle}
\end{equation}
The constant-spectrum case is handled separately. Degree selection evaluates
logarithms without first forming a potentially underflowing ratio of a very
small tolerance to a tail factor.

The default compact degree budget is 1023, with a default estimated working
memory budget of 512 MiB. A $1000\times1000$ diagonal SPD input with condition
number 10000 requires degree 1842 at the $10^{-12}$ truncation target and is
explicitly rejected by the degree guard. This is a configured resource limit,
not a proof that the underlying approximation is impossible. Raising a guard
does not remove conditioning or finite-precision limitations.

\subsection{Exact success projection in a structured statevector}
Let the encoded state after SELECT have amplitudes $A_{k,f,j}$, with selector
$k$, flag $f$, and spectral data index $j$. The zero-selector amplitude after
PREP$^\dagger$ is exactly
\begin{equation}
 B_{f,j}=\sum_k w_k^*A_{k,f,j}.
 \label{eq:projection}
\end{equation}
The simulator evaluates all controlled rotation gates on the full
$K\times2\times P$ state, where $K=2^s$, and obtains the successful data
amplitudes from $B_{0,j}$. Other sectors retain a lazy selector basis change.
State norm is measured before that norm-preserving change. This eliminates
the need to build and multiply a dense $K\times K$ selector matrix during
normal large runs, without evaluating the reciprocal polynomial as a
classical substitute for the simulated gates.

For small comparisons the code explicitly constructs PREP$^\dagger$ and
returns every post-uncomputation selector amplitude. Tests compare this full
state with the lazy success projection and with Qiskit native-statevector
execution. The final $V$ basis change is evaluated on the success sector;
the other sectors retain a lazy output basis. Thus the large simulator is
an exact structured ideal circuit simulation, not native-gate Aer execution
at dimension 1000 and not a noisy hardware simulation.

The stored state occupies $16(2KP)$ bytes in complex128 precision. The
preflight working-memory estimate is six such states plus $64P^2$ bytes for
dense encoding and temporary workspace. This is a conservative planning
estimate, not an operating-system allocation cap or a model of native
compiler memory. The largest recorded compact case has 21 qubits, a 32 MiB
statevector and estimated 256 MiB working memory.

\subsection{Complete-operator audit}
Three input vectors alone do not certify a matrix inverse. To audit every
direction, a balanced input is prepared in the spectral basis so that all
$N$ valid coordinates are nonzero. SELECT is diagonal in that data basis;
dividing simulated success amplitudes by the known actual input amplitudes
exposes every diagonal response of the implemented block. The audit forms
$\widetilde X=V\operatorname{diag}(g_j)U^T$ and evaluates
\begin{equation}
 \norm{M\widetilde X-I}_2,\quad
 \frac{\norm{M\widetilde X-I}_F}{\sqrt N},\quad
 \frac{\norm{\widetilde X-X_{\rm ref}}_F}{\norm{X_{\rm ref}}_F}.
 \label{eq:audit}
\end{equation}
The independent $X_{\rm ref}$ is computed only after circuit construction
and simulation. No response is fed back to modify the circuit. The first
metric is an observed induced-norm bound for all inputs of that tested
matrix, including floating-point basis-change error.

This audit exploits simulator amplitudes and the documented diagonal
structure. It is not tomography, device-derived fidelity, or a scalable
physical readout of $N^2$ classical inverse entries. Hardware verification
must use measurements and account for the depth and uncertainty of its
verification circuit.

\section{Experimental Design and Reproducibility}
\subsection{Fixtures and thresholds}
The dimension study uses real matrices of orders 10, 100 and 1000 in four
families: dense SPD, tridiagonal SPD, dense nonsymmetric, and signed
indefinite. Fixed-seed small perturbations are added to diagonals spanning
1.5 to 2.5; tridiagonal off-diagonal entries are 0.08. Exact generators and
seed rules are distributed with the code \cite{qcenncode}. These are
well-conditioned fixtures, not samples from all nonsingular matrices.

The precision study augments the four order-1000 fixtures with dense
controlled-spectrum matrices. Fixed-seed QR factors provide orthogonal
left/right bases; singular values form a geometric sequence from 1 to
$\kappa$. For each $\kappa\in\{10,100,1000\}$, we generate SPD,
nonsymmetric, and symmetric indefinite examples. Alternating eigenvalue
signs define the indefinite family while preserving nonsingularity.
Basis, uniform and random inputs are tested in addition to the complete
operator audit. A like-for-like baseline repeats those inputs with a
$10^{-6}$ polynomial target.

The tighter target is $\delta_d\leq10^{-12}$. The distinct numerical
acceptance gate is $10^{-10}$ on measured operator residual, sampled
relative residual and relative full-inverse error. The analytic tail bound
does not cover SVD error, angle rounding or floating-point arithmetic.
State norm and padding leakage are checked separately. Consequently, a
passing polynomial formula never suffices to pass the numerical experiment.

The theory regression spans orders 2 and 4, three matrix types, condition
numbers 2, 10 and 100, and nine degrees from 1 to 511. Its
\TheoryCaseCount{} cases test the operator, per-input, normalized-state,
and success inequalities with explicit $10^{-10}$ floating-point slack.
An inaccurate low-degree approximation can correctly satisfy a loose
analytic bound; such agreement is not labeled target-accuracy success.

\subsection{Execution and provenance}
Qiskit supplies circuit construction and native execution interfaces
\cite{qiskit2024}. The recorded run uses Python \PythonVersion{}, NumPy
\VersionNumpy{}, SciPy \VersionScipy{}, Qiskit \VersionQiskit{}, and
Qiskit Aer \VersionQiskitaer{}. OpenBLAS, OpenMP and MKL thread controls are
set to one. A single Compose service runs static checks, \TestCount{} unit
tests, A1, coupled A2--A4, block-hybrid scaling, coherent Neumann baselines,
compact precision, theory regression and native-noise preparation.

Each stage runs in a fresh process, retaining its command, return code,
stdout/stderr, elapsed wall time, CPU time and peak child resident memory.
The evidence includes JUnit XML, raw-count audits, QPY circuits, JSON
metrics, package versions, source commit and source-tree SHA-256. Independent
campaign failures are retained rather than silently stopping the remaining
tests. Software correctness and synthetic noisy-quality decisions have
separate fields and cannot be inferred from one another.

The Python dependency closure is pinned. The container's base image can be
selected by immutable digest for exact replay; a floating default base tag
is not a promise of bitwise reproducibility across future builds. The
manifest and immutable commit identify the actual experiment. The
publication builder derives tables and figures from those run outputs
instead of copying manually selected numbers into the manuscript.

\section{Results}
\subsection{Approximation and order-1000 precision}
Fig.~\ref{fig:convergence} compares the simulated operator residual with
the truncation bound. Increasing degree reduces approximation error until
floating-point effects dominate. The conservative tail bound varies with
conditioning; degree alone is not an accuracy certificate without an
interval specification. The regression includes initially poor approximants
to make this dependency observable.

\begin{figure*}[t]
\centering\includegraphics[width=0.85\textwidth]{figures/convergence.pdf}
\caption{Structured gate-simulation residual and analytic truncation bound
for nonsymmetric order-4 matrices with controlled condition numbers.
Curves below the numerical floor are clipped at $10^{-16}$ for display.
The clipping is visual only; raw values remain in the JSON evidence.}
\label{fig:convergence}
\end{figure*}

Table~\ref{tab:precision} reports all \PrecisionCaseCount{} order-1000
operators. The largest measured residual is \WorstOperatorResidual{};
the largest relative Frobenius inverse error is \WorstInverseError{}.
The original well-conditioned families use degrees 13--16 and 15--16
total qubits. Controlled condition numbers 10, 100 and 1000 require degrees
46, 161 and 546 and use 17, 19 and 21 qubits, respectively.

\begin{table*}[t]
\caption{Complete simulated inverse-operator audit for $N=1000$.
The final column is the minimum success probability over normalized valid inputs.}
\label{tab:precision}\centering\small
\begin{tabular}{lrrrrrr}\toprule
Family & $\kappa$ & $d$ & Qubits & $\norm{M\widetilde X-I}_2$ &
Relative inverse error $F$ & Minimum success\\\midrule
\input generated/precision_rows.tex 
\bottomrule\end{tabular}
\end{table*}

On the four earlier families, the worst residual over the same three
sampled inputs is reduced by approximately $2.3\times10^5$ to
$1.6\times10^6$ relative to the looser $10^{-6}$-target baseline.
This comparison changes both degree and numerical precision safeguards;
it is not a measurement of hardware accuracy or an improvement over a
classical direct solver. Fig.~\ref{fig:precision} shows that the
high-condition cases remain within the separate numerical gate. Numerical
residuals need not decrease monotonically with $\kappa$ because different
adapted degrees and floating-point effects are involved.

\begin{figure*}[t]
\centering\includegraphics[width=0.85\textwidth]{figures/precision.pdf}
\caption{Order-1000 controlled-spectrum precision cases. The truncation target
and numerical acceptance limit are different quantities. Observed values
come from the reconstructed simulated success operator.}
\label{fig:precision}
\end{figure*}

The selector-memory improvement makes these higher degrees practical in
the structured simulator. It does not reduce the dense encoding to an
efficient oracle. At $\kappa=1000$, the worst-input success is about one
accepted output per million independent attempts. Fig.~\ref{fig:resources}
juxtaposes the modest raw-qubit growth with that sampling penalty.

\begin{figure*}[t]
\centering\includegraphics[width=0.85\textwidth]{figures/resources.pdf}
\caption{Raw qubits and minimum success for the order-1000 nonsymmetric
controlled-spectrum family at the tighter target. The near-$\kappa^{-2}$
sampling cost remains even though the selector register grows slowly.}
\label{fig:resources}
\end{figure*}

\subsection{Native compilation and synthetic noise}
Table~\ref{tab:native} lists compact action circuits compiled to
$\{R_z,\mathrm{SX},X,\mathrm{CX}\}$ on a bidirectional active-width line.
Final routing permutations are restored before an ideal statevector
residual is calculated. This layout is an explicit generic assumption,
not a calibrated device target. The native campaign uses a $10^{-6}$
polynomial target to control gate cost; its circuits are not the
high-precision order-1000 circuits in Table~\ref{tab:precision}.

\begin{table}[t]\caption{Generic-line native action circuits, target $10^{-6}$.}
\label{tab:native}\centering\footnotesize
\begin{tabular}{rrrrr}\toprule
$N$ & Qubits & CX & Depth & Ideal residual\\\midrule
\input generated/hardware_rows.tex 
\bottomrule\end{tabular}
\end{table}

The independent verifier appends a reference-state preparation inverse
after the solver. Ancilla success is counted first; among accepted shots,
an all-zero reference return is measured. Its conditional fraction includes
reference-verifier and readout errors. It is not a tomography-derived
fidelity and is not used to fit or correct the inverse action.

Two synthetic profiles use depolarizing probabilities of $10^{-5}$/$10^{-4}$
on SX/X and $10^{-4}$/$10^{-3}$ on CX, with symmetric readout flips of
0.005/0.02, respectively. Virtual $R_z$ rotations are noiseless in this
model. Orders 2, 4 and 10 use density-matrix channel simulation with 32768
measurement shots. Order 100 uses only the optimistic profile, simulated
as 64 noisy statevector trajectories to bound runtime.

The provisional quality gate requires at least 20 accepted shots and a
95\% Wilson lower confidence bound of at least 0.9 on conditional return
\cite{wilson1927}. It is a project decision rule, not a universal hardware
standard. Table~\ref{tab:noise} and Fig.~\ref{fig:noise} retain all outcomes.
The order-10 circuit passes the optimistic profile but fails the stress
profile. The dense order-100 circuit fails the optimistic profile despite
passing ideal native correctness. Its small trajectory sample produces a
wide interval, so the observation is a bounded negative indication rather
than a precise device-performance forecast.

\begin{table*}[t]\caption{Synthetic native-gate verifier results. Gate decisions
use the confidence bound and minimum accepted-shot count, not the point estimate alone.}
\label{tab:noise}\centering\small
\begin{tabular}{rlrrrll}\toprule
$N$ & Profile & Shots & Accepted & Conditional return & 95\% interval & Gate\\\midrule
\input generated/noise_rows.tex 
\bottomrule\end{tabular}
\end{table*}

\begin{figure*}[t]\centering
\includegraphics[width=0.85\textwidth]{figures/noise.pdf}
\caption{Observed conditional reference return with 95\% Wilson intervals
under synthetic native-gate noise. The verifier contributes noise. The
order-100 optimistic point uses only 64 trajectories; no stress point at
that order is implied.}
\label{fig:noise}\end{figure*}

\section{Resource Interpretation and Hardware Study}
\subsection{Costs not captured by the qubit count}
Dense eigendecomposition/SVD requires conventional cubic preprocessing and
quadratic storage. An unstructured real orthogonal basis has
$N(N-1)/2$ continuous parameters, so representing it with $\log_2N$ data
qubits does not make generic synthesis logarithmic. A dense padded basis
change typically has quadratic gate cost in $P$. Multiplexed walk gadgets
have a linear number of listed spectral angles in $P$ per selector bit.
The binary-power construction benefits from known classical angles and
cannot be compared directly to black-box queries without charging for
that information and its precision.

Input loading, accepted-output sampling, verification, and classical
readout are additional costs. An entire inverse contains $N^2$ classical
entries. Reporting a small inverse-action register never establishes
efficient measured extraction of that matrix. A useful hardware study
should therefore select specific output tasks and structured input
families, rather than equate qubit capacity with full inversion throughput.

\begin{proposition}[Sensitivity to coherent implementation error]
Let ideal and implemented unitaries differ in operator norm by at most
$\eta$. For a fixed normalized input with ideal postselection probability
$p>0$, if $\eta<\sqrt p$, the normalized accepted output vectors differ
by at most $2\eta/\sqrt p$.
\end{proposition}
\begin{proof}
An orthogonal success projection is a contraction, so the two unnormalized
successful vectors differ by at most $\eta$. Their ideal norm is $\sqrt p$;
the assumption ensures the implemented vector is nonzero. Apply the
normalization inequality used in~\eqref{eq:statebound}.
\end{proof}
For products of unitary gates, a telescoping expansion gives
$\eta\leq\sum_\ell\eta_\ell$. This illustrates why deep circuits and small
success probabilities are unfavorable together. It is not a model of
general stochastic noise, mixed-state postselection, or readout error;
the synthetic experiments are reported separately for that reason.

\subsection{A staged physical validation protocol}
Physical validation is proposed first for orders 2 and 4, followed by
order 10 only if the smaller cases meet prespecified measurement gates.
For each target backend, the unfused source circuit is regenerated and
transpiled against its actual native gate set, coupling graph and current
calibration. Previously exported generic-line QPY files are evidence of
compilation, not device-specific executable certification.

The first stage should measure prepared-input and reference-verifier
controls, then the inverse-action circuits for basis and superposition
inputs. It should retain raw counts, calibration timestamps, routing,
native operation counts, depth, shots, job IDs, failure records and
confidence intervals. Small-register signed-state tomography is needed
where a full inverse residual is claimed; a reference echo alone does
not supply it. A paired calibration-derived emulator should precede
submission and be compared with the measured result without substituting
for it.

Qubit capacity, useful precision, and useful acceptance rate are separate
criteria. Dense orders 100 and 1000 should not be promised as successful
physical experiments on the basis of these ideal statevectors. The current
order-100 noisy failure motivates reducing dense encoding depth, studying
structured inputs and analyzing appropriate preconditioning before scaling
physical submissions. Any preconditioner must be charged to preprocessing
and loading; it must not conceal a classical solution inside the input.

The accompanying report is organized for the EuroHPC Quantum Access Pilot
technical-feasibility process \cite{eurohpc2026}. It distinguishes applicant
allocation information from measured resource use and identifies the
hosting-system evidence that remains to be collected. Neither this paper
nor that report claims a physical allocation decision or execution on a
specific allocated EuroHPC node.

\section{Discussion and Conclusion}
The specialized certification and implementation show that finite real
nonsingular matrices admit coherent approximate inverse action in the
specified spectral access model. The demonstration is stronger than a
few successful right-hand sides because the large-instance audit measures
the induced residual of each full simulated success operator. It is also
more limited than a scalable quantum inversion claim: dense classical
encoding is required, the quantum output is postselected, and full measured
matrix readout is absent.

In the recorded experiment, \PrecisionCaseCount{} order-1000 operators
reach a largest residual of \WorstOperatorResidual{} and a largest relative
Frobenius inverse error of \WorstInverseError{}. The complete suite includes
\TestCount{} unit tests and explicit bound regressions. The raw register
count remains 15--21 qubits for the high-precision cases, while worst-input
success falls to approximately $10^{-6}$ at condition number 1000.
Native synthetic-noise failures are an essential part of the result,
demonstrating that ideal numerical precision is insufficient evidence
for hardware utility.

The next research step is calibrated small-instance physical validation
with reproducible measurement evidence. Longer-term work should replace
dense spectral access with structured block encodings, evaluate the cost
of amplification or preconditioning, and choose observables that can be
estimated without full inverse readout. The present artifact provides a
testable baseline for those comparisons rather than removing their
resource constraints.

\appendix[Reproduction and Evidence Integrity]
The default command is
\begin{quote}\small\texttt{docker compose up --build}\\
\texttt{--abort-on-container-exit}\\
\texttt{--exit-code-from validation}.
\end{quote}
Outputs are written to the mounted artifact directory. A fresh output
directory is required for a new run to prevent silent mixing of stale
results. The same Python suite can be executed on a system where Docker
is not available, including an HPC environment using its supported
container conversion or a managed Python environment. Actual site
software and scheduler policy must be checked before deployment.

This manuscript was prepared for reporting on \ReportDate{} from a run whose
execution-environment clock reported \RunDate{}, source commit
\texttt{\SourceCommit}, with execution label \texttt{\RunContext}.
The source-tree fingerprint and full commands are retained in the
artifact, including an explicit dirty-tree indicator for local development
runs. The measured benchmark portion used \SuiteWallSeconds{} wall seconds,
\SuiteCPUSeconds{} CPU seconds and a maximum per-stage child RSS of
\SuitePeakMiB{} MiB. These figures are local process measurements, not
EuroHPC node-hour accounting and not physical QPU runtime. Document
rendering is excluded from these benchmark totals.
The raw process clock is not independently synchronized; scheduler-attested
execution dates are not inferred from it.

\section*{Acknowledgment and AI Assistance Disclosure}
OpenAI ChatGPT/Codex assisted in drafting and revising Sections I--IX and
the appendix, developing the validation/publication scripts, and preparing
the mathematical exposition. Figures are deterministic plots of archived
numerical outputs, not generated illustrative images. This disclosure
identifies the use of AI assistance in accordance with IEEE guidance
\cite{ieeeai}. Author affiliations and final venue-specific submission
metadata are not inferred from the software repository.

\bibliographystyle{IEEEtran}
\bibliography{references}
\end{document}
